% =========================================================================
% บทที่ 7: ฟิสิกส์ของแข็งและวัสดุแม่เหล็กขั้นสูง
% =========================================================================

\chapter{ฟิสิกส์ของแข็งและวัสดุแม่เหล็กขั้นสูง}
\label{ch:chapter07}

\noindent{\large\bfseries\color{physicsblue} โครงสร้างผลึก ทฤษฎีแถบพลังงาน อันตรกิริยาแลกเปลี่ยน และการจำลองเชิงทฤษฎี LSDA+U}

\vspace{0.4cm}

\begin{equationsbox}{แบบจำลองแม่เหล็กและฮามิลโทเนียนของไฮเซนเบิร์ก}
\begin{align}
\hat{H}_{\text{Heisenberg}} &= -2\sum_{i < j} J_{ij} \mathbf{S}_i \cdot \mathbf{S}_j, \quad E^{\text{LSDA}+U} = E^{\text{LSDA}} + \frac{U - J}{2}\sum_{\sigma}\left[\text{Tr}(\mathbf{n}^\sigma) - \text{Tr}(\mathbf{n}^\sigma \mathbf{n}^\sigma)\right]
\end{align}
\end{equationsbox}

\section{ทฤษฎีและกรอบแนวคิดสำคัญ}

ฟิสิกส์ของแข็ง (Solid-State Physics) ศึกษาพฤติกรรมของอิเล็กตรอนในโครงตาข่ายผลึกที่มีความสมมาตรแบบเลื่อนขนาน (Translational Symmetry) ตามทฤษฎีของโบลช (Bloch's Theorem) ฟังก์ชันคลื่นของอิเล็กตรอนในผลึกจะอยู่ในรูปของคลื่นระนาบปรับสัญญาณด้วยฟังก์ชันรายคาบของโครงตาข่าย นำไปสู่การจำแนกสารเป็นตัวนำ กึ่งตัวนำ และฉนวน ตามโครงสร้างแถบพลังงาน (Energy Band Structure)

อย่างไรก็ตาม ในระบบที่มีความสัมพันธ์ของอิเล็กตรอนอย่างเข้มข้น (Strongly Correlated Electron Systems) เช่น ออกไซด์ของโลหะทรานสิชัน มอทท์-ฮับบาร์ด (Mott-Hubbard Insulators) อย่าง $\text{MnO}$ และ $\text{NiO}$ ทฤษฎีฟังก์ชันความหนาแน่นดั้งเดิม (Standard DFT/LSDA) มักล้มเหลวโดยทำนายว่าสารเหล่านี้เป็นตัวนำกึ่งโลหะ ซึ่งขัดแย้งกับผลการทดลองจริงที่เป็นฉนวนแม่เหล็กแอนติเฟร์โร (Antiferromagnetic Insulators) ที่มีช่องว่างแถบพลังงานกว้าง การแก้ไขปัญหานี้ต้องใช้ระเบียบวิธี $\text{LSDA}+U$ ซึ่งเพิ่มพจน์พลังงานผลักคูลอมบ์ในออร์บิทัล $d$ ที่มีค่า $U$ (On-site Coulomb Interaction) และอันตรกิริยาแลกเปลี่ยนฮุนด์ $J$ (Hund's Exchange Interaction)

\section{ตัวอย่างการคำนวณตามกรอบวิเคราะห์ P-S-C-T}

\begin{psctexample}{7.1}{การประเมินอันตรกิริยาแลกเปลี่ยนแม่เหล็ก J1 และ J2 ในสารประกอบ MnO และ NiO}
\textbf{1. ภาพและสถานการณ์ทางกายภาพ (Picture):}\\\\ โครงสร้างแม่เหล็กแบบ AF II ของ MnO บนโครงตาข่าย fcc มีอันตรกิริยาแลกเปลี่ยนระหว่างเพื่อนบ้านลำดับหนึ่ง ($J_1$) และลำดับสอง ($J_2$)

\vspace{4pt}
\textbf{2. กลยุทธ์และแบบจำลองทางฟิสิกส์ (Strategy):}\\\\ เปรียบเทียบพลังงานรวมจากการจัดเรียงสปินแบบเฟร์โร (FM) และแบบแอนติเฟร์โร (AFM) เพื่อสกัดค่า $J_1$ และ $J_2$

\vspace{4pt}
\textbf{3. ขั้นตอนการคำนวณเชิงตัวเลข (Calculation):}\\\\ พลังงานต่อไอออนแม่เหล็กตามแบบจำลองไฮเซนเบิร์ก:
\begin{align}
E_{\text{FM}} &= E_0 - 6 J_1 S^2 - 3 J_2 S^2 \\
E_{\text{AF II}} &= E_0 + 3 J_2 S^2
\end{align}
จากผลงานวิจัยของ ผศ.ดร.ชีวะ ทัศนา ด้วยระเบียบวิธี FP-LAPW $\text{LSDA}+U$:
ผลต่างพลังงานระหว่างสถานะแม่เหล็กนำไปสู่การคำนวณค่า $J_2 \approx -0.45\text{ meV}$ ซึ่งมีเครื่องหมายลบ บ่งชี้ว่าอันตรกิริยาแลกเปลี่ยนระหว่างอะตอมเพื่อนบ้านลำดับสองผ่านสะพานออกซิเจน $\text{Mn}-\text{O}-\text{Mn}$ (Superexchange) เป็นอันตรกิริยาหลักที่ทำให้เกิดโครงสร้างแม่เหล็กแบบแอนติเฟร์โรที่เสถียร

\vspace{4pt}
\textbf{4. การทดสอบความสมเหตุสมผล (Test):}\\\\ ค่าอุณหภูมิเนเอล (Néel Temperature: $T_N$) ที่คำนวณจากทฤษฎีสนามเฉลี่ย (Mean-Field Theory) โดยใช้ค่า $J$ ที่สกัดได้ สอดคล้องกับค่าจากการทดลอง $T_N \approx 118\text{ K}$ ใน $\text{MnO}$
\end{psctexample}

\section{การเชื่อมโยงสู่การวิจัยฟิสิกส์ร่วมสมัย}

\begin{researchbox}{ข้อมูลเชิงลึกจากงานวิจัย}
ผลงานวิจัยของ ผศ.ดร.ชีวะ ทัศนา ที่ตีพิมพ์ในวารสารระดับนานาชาติ ได้แสดงให้เห็นว่าการเลือกค่าพารามิเตอร์ $U$ และ $J$ ที่เหมาะสมในระเบียบวิธี $\text{LSDA}+U$ ไม่เพียงแต่เปิดช่องว่างพลังงาน Band Gap ใน $\text{MnO}$ และ $\text{NiO}$ ได้อย่างถูกต้องเท่านั้น แต่ยังทำนายค่าโมเมนต์แม่เหล็กเฉพาะที่ (Local Magnetic Moments) ได้ตรงกับผลการทดลองการกระเจิงนิวตรอน (Neutron Scattering)
\end{researchbox}

\section{การฝึกแก้โจทย์ปัญหาเชิงระบบตามกรอบ C-C-A-F}

\begin{ccafsolution}{การวิเคราะห์โครงสร้างแถบพลังงานตามแบบจำลองโครนิก-เพนนี}
\textbf{ขั้นที่ 1: การสร้างมโนทัศน์ (Conceptualize):}\\\\ อิเล็กตรอนเคลื่อนที่ผ่านแนวกำแพงศักย์รูปฟังก์ชันเดลตาแบบรายคาบ $V(x) = V_0 b \sum_n \delta(x - na)$

\vspace{4pt}
\textbf{ขั้นที่ 2: การจำแนกประเภทโจทย์ (Categorize):}\\\\ การแก้สมการโบลชและการเกิดช่องว่างแถบพลังงาน (Kronig-Penney Model \& Band Gaps)

\vspace{4pt}
\textbf{ขั้นที่ 3: การวิเคราะห์เชิงคณิตศาสตร์ (Analyze):}\\\\ สมการเงื่อนไขพลังงานของโครนิก-เพนนี:
\begin{equation}
P \frac{\sin(\alpha a)}{\alpha a} + \cos(\alpha a) = \cos(ka), \quad P = \frac{m V_0 b a}{\hbar^2}
\end{equation}
เนื่องจากฟังก์ชัน $\cos(ka)$ มีค่าอยู่ระหว่าง $[-1, 1]$ ดังนั้น ค่าของ $\alpha a$ ที่ทำให้ฟังก์ชันทางซ้ายมือมีค่าสัมบูรณ์มากกว่า 1 จะเป็นช่วงพลังงานที่อิเล็กตรอนไม่สามารถอยู่ได้ เกิดเป็นช่องว่างแถบพลังงานต้องห้าม (Forbidden Energy Gap)

\vspace{4pt}
\textbf{ขั้นที่ 4: การสรุปและประเมินผล (Finalize):}\\\\ แบบจำลองโครนิก-เพนนียืนยันว่าช่องว่างแถบพลังงานเป็นผลตามธรรมชาติของการแทรกสอดคลื่นอิเล็กตรอนในโครงตาข่ายผลึกรายคาบ
\end{ccafsolution}

\section{แบบฝึกหัดทบทวนและโจทย์ท้าทายประจำบท}

\begin{enumerate}[leftmargin=1.5cm, label={\textbf{\thechapter.\arabic*}}, itemsep=8pt]
    \item \textbf{โจทย์เชิงแนวคิด (Conceptual):} จงอธิบายความสำคัญทางกายภาพของกฎและสมการหลักในบทเรียนนี้ พร้อมยกตัวอย่างสถานการณ์จริงในธรรมชาติ
    \item \textbf{โจทย์การประมาณค่าเชิงฟิสิกส์ (Estimation):} จงใช้การวิเคราะห์เชิงมิติและอันดับของขนาด (Order of Magnitude) ในการประมาณค่าปริมาณสำคัญที่เกี่ยวข้อง
    \item \textbf{โจทย์วิจัยขั้นสูง (Research Challenge):} จงเขียนสคริปต์ Python จำลองพฤติกรรมของระบบตามสมการที่ได้ศึกษา พร้อมพลอตกราฟเปรียบเทียบผลลัพธ์เชิงทฤษฎี
\end{enumerate}

